\section{Further research questions and proposed next steps}
\label{sec:future}

The exact measure separates several questions that were difficult to
distinguish from the original recurrence. The following are concrete
extensions, with their unresolved parts stated explicitly.

\begin{enumerate}[leftmargin=*]
\item \textbf{Complete the priority comparison.}
Obtain the full Ford--Lau solution and inspect subsequent work on this
specific recurrence, including citations of the original problem.
Determine precisely which of the density formula, complete
monotonicity, all-orders coefficients, and product correction have
appeared before. This is the first step before describing the result
as a first publication.

\item \textbf{Optimal truncation of the inverse-logarithmic expansion.}
The coefficient generator is an entire reciprocal-gamma expression,
but the asymptotic coefficients also contain $(k+1)!$. Determine their
large-order behavior, the optimal truncation index as a function of
$\log n$, and the least attainable remainder for this single series.
The exact $F_a$ integrals provide a reference against which a proposed
optimal-truncation theorem can be tested. Fixed-order error estimates
in this paper do not answer this question.

\item \textbf{Uniformity when the algebraic order grows.}
Theorem~\ref{thm:main-sectors} and its general version hold for each
fixed truncation order. Quantify the dependence of the remainder
constant on that order, and determine a range $R=R(n)\to\infty$.
The singularities at distance $2\pi$ in the endpoint correction
$P$ are natural inputs, but the logarithmic denominator must be
controlled simultaneously.

\item \textbf{Interval-certified product constants and errors.}
Implement directed-rounding enclosures for $\rho$, $\gamma$, and
the positive integral. Combine them with
\eqref{eq:C-enclosure} or \eqref{eq:alternating-tail} to certify
the infinite product $C_*$ to an arbitrary requested number of digits.
An implementation should avoid subtraction of two exponentially large
terms. The current scripts demonstrate numerical agreement; they do
not provide these interval guarantees.

\item \textbf{Growing finite-difference order.}
The exact formula for $D_{n,r}$ remains valid for every $r$, whereas
\eqref{eq:difference-asymptotic} assumes fixed $r$. When $r/n$ stays
away from zero and infinity, its integral is controlled by an interior
beta-distribution saddle. Derive a uniform theorem joining that
interior regime to the logarithmic endpoint regime $r=o(n)$.
Such a theorem would describe the full array of difference inequalities,
not just each fixed column.

\item \textbf{General logarithmic renewal kernels.}
Suppose $a_{j+1}-a_j=\int_0^1t^{j-1}h(t)\,dt$ for a positive
density $h$, and $1+a_1\ge\int_0^1h(t)/t\,dt$.
Theorem~\ref{thm:general-kernel} gives an all-orders logarithmic
expansion under explicit regularity assumptions, with a possible atom
at zero. Determine minimal endpoint regularity for that conclusion,
and characterize precisely which deformations preserve the special
cancellation of Euler's constant in \eqref{eq:main-coeff}.
If $h(1)>0$ and a suitable logarithmic endpoint limit exists, the
same calculation predicts leading constant $1/h(1)$.

\item \textbf{Stability under perturbing initial coefficients.}
A finite change of the recurrence coefficients leaves their
logarithmic tail unchanged but may create additional poles.
Classify the perturbations for which the positive remainder measure
survives, and those for which oscillating exponential corrections
appear. The singularity at one should still determine a universal
logarithmic component after all relevant poles are removed; that
assertion requires a separate global zero analysis.

\item \textbf{Arithmetic error propagation.}
Insert the proved product correction into a fully audited
totient-volume argument and track every competing approximation.
Determine whether any arithmetic regime exposes the new
$\rho^m/(m^2\log^2m)$ scale rather than hiding it under a larger
sieve or boundary error. This paper proves the continuous-scale
correction and leaves that arithmetic comparison open.

\item \textbf{Lean formalization in stages.}
Start with the convolution recurrence, the integral formula for its
increments, and the algebraic relation between $G$ and $J$.
Next formalize the positive bounded operator and its scalar spectral
measure, including endpoint atoms. The exact moment inequalities
then follow with little asymptotic machinery. A final stage can
formalize the uniform Laplace estimates and coefficient extraction.
No declaration in the supplied package is represented as a
kernel-checked proof.
\end{enumerate}

\section{Integration and reproducibility record}
\label{sec:integration}

The package is self-contained for compilation and for the supplied
numerical experiments. It contains the complete LaTeX source, a
compiled PDF, vector figures and generated tables, Python programs,
CSV/JSON data, a build file, and a claim-status ledger. No original
repository manuscript or third-party PDF is redistributed.

The inspected repository snapshots were
\begin{center}\small
\begin{tabular}{@{}lp{0.71\linewidth}@{}}
\toprule
Repository & Commit\\
\midrule
\texttt{openai/math}&\path{adc7f1241b42e322a6451854ab7e4b4c146bf78a}\\
\texttt{ProveIt}&\path{6d78805163d6b5b92d27f8f41555254c49fe0963}\\
\bottomrule
\end{tabular}
\end{center}
The first snapshot supplies the recurrence and product-scale anchor;
the second records the repository scope at the start of the study.
Their full paths and the relevant source blob identifier appear in
\texttt{provenance.json}. The inspection was targeted, and did not
constitute a theorem-by-theorem audit of either repository.

An appropriate proposed destination is
\path{Analysis/RenewalTheory/Research/Ford_Renewal_Logarithmic_Remainder/}.
This is an integration suggestion; no existing repository content has
been changed or published by this package.

The mathematical dependency is short: a positive operator gives the
exact residual measure; elementary endpoint estimates give its sharp
size and explicit bounds; a convergent polylogarithm identity gives
the complete expansions; positive integral identities transport them
to products and renewal accuracy. The numerical programs check each
description against an independently computed representation. Their
purpose is reproducibility and error detection, while the proofs above
establish the infinite statements.
